\section{Key Trends and Patterns Across Results}

\subsection{The Field Is Optimizing the Smaller of Two Effects}
The central pattern is a mismatch between where effort is spent and where performance is
determined. Architectural work is abundant: the corpus contains transformers, capsule
networks, autoencoders, generative augmentation and federated schemes, and each reports an
improvement over its chosen baseline. Yet on the one dataset where these architectures can be
compared, they occupy a two-point band that a $k$-nearest-neighbour classifier also occupies
\cite{r14,r50}. Meanwhile the factors that move performance by an order of magnitude more,
whether the model is personalized, whether evaluation is subject independent, whether data
were collected in the field, are treated as reporting details and are usually not reported at
all. A median fusion gain of $4.67$ pp against a median evaluation gap of $9.14$ pp, and
against a personalization gap of $38.03$ pp, quantifies that mismatch.

\subsection{Fusion Is Real, Bounded and Baseline-Dependent}
The four protocol-matched contrasts are unanimously positive, so the case for multimodal
sensing is not overturned. It is bounded. Fusion buys most where the unimodal baseline is
weak and the added channel is physiologically distinct, as in EEG supplemented by GSR at
$+14.30$ pp \cite{r75}, and very little where the baseline is already strong, as in GSR
supplemented by EEG at $+1.00$ pp of AUROC \cite{r37}. Adding a camera channel to a
physiological wristband returned $+3.23$ pp in the one corroborated office deployment
\cite{r26}. For system designers the operational reading is that a second sensor must be
justified against the specific unimodal baseline it will actually replace, together with its
cost in power, cost, form factor and failure modes, and not against a generic claim that
fusion improves performance.

\subsection{Personalization Is the Dominant Lever, and the Main Threat to Validity}
Personalization accounts for the largest effect in the corpus, and simultaneously for the
most serious threat to the interpretability of the literature. Because subject-dependent
splits allow the same participant to appear in training and test, they can be reported
without any deliberate misrepresentation while producing figures that a subject-independent
evaluation would not support. The one study in the corpus that reports both conditions on
identical data shows accuracy falling from $95.06\%$ to $67.65\%$ and F1 from $91.71$ to
$43.05$ \cite{r33}. Since only 8 of 60 studies state their protocol recoverably, the scale of
this effect across the wider literature cannot currently be estimated, which is itself the
finding.

\subsection{Laboratory Performance Does Not Forecast Field Performance}
Two independent within-study contrasts show the same direction and a similar magnitude:
$13.00$ pp of F1 lost between laboratory and free-living conditions \cite{r13}, and
$14.56$ pp of AUROC lost between WESAD and a real-world student cohort \cite{r49}. Transfer
between two curated benchmarks, by contrast, cost only $4.00$ pp \cite{r55,r61}. The gap
between these two figures is the crux: a study that demonstrates cross-dataset transfer
between laboratory corpora has not demonstrated the property that deployment requires. The
structure of the corpus makes this consequential, since 31 of 60 studies are laboratory or
benchmark-only and just 9 involve free-living deployment.

\subsection{Benchmark Saturation Has Made Architectural Claims Unfalsifiable}
When a benchmark compresses a decade of architectural development into a two-point band, it
cannot adjudicate between methods, and improvements reported on it are no longer informative
about which method is better. WESAD has reached that state for subject-dependent and LOSO
evaluation. It has not reached it for subject-independent generalization, where the single
available record sits at $67.65\%$ \cite{r33}. The productive response is therefore not to
abandon WESAD but to change the protocol reported on it: subject-independent generalization
on WESAD remains a discriminating and unsolved task.

\subsection{Reporting Practice Is the Rate-Limiting Step}
The reporting audit reframes every preceding trend. A validation protocol recoverable for
$13.3\%$ of studies and a variance estimate for $5.0\%$ means that most reported accuracies
in this area cannot be placed on a common scale, cannot be meta-analyzed defensibly, and
cannot be used to specify a system. This is why the present review restricts itself to
within-study contrasts and to 28 corroborated studies rather than pooling 60 headline
numbers: the pooled figure would be larger, more impressive and not interpretable.

\section{Discussion}

\subsection{Summary of Findings}
This review set out to establish whether the performance benefit attributed to multimodal
fusion in wearable and camera-based stress and attention monitoring survives a
confound-controlled estimate, and how it compares with the influence of evaluation design.
Restricting the synthesis to within-study contrasts and to independently corroborated values
yields three findings.

First, fusion helps, by a median of $4.67$ pp across four protocol-matched contrasts, with
the gain concentrated where the unimodal baseline is weak. Second, evaluation choices move
reported performance by a median of $9.14$ pp, roughly twice as much, and personalization
alone by $38.03$ pp, up to $10.4$ times the median fusion gain. Third, the benchmark on which
most architectural claims rest has saturated for the protocols commonly used on it, while
remaining unsolved for the subject-independent protocol that deployment requires, and the
protocol actually used is unrecoverable for the large majority of studies.

Taken together these support a single conclusion with direct implications for BSPC's
translational remit: the binding constraint on deployable stress and attention monitoring is
not the sensing configuration or the network architecture but the evaluation regime under
which systems are developed and reported. A system optimized against a subject-dependent
score on a saturated benchmark is being optimized against a target that does not predict
field behavior.

Three concrete recommendations follow. Report subject-independent performance as the
primary result: personalized and subject-dependent figures may be reported alongside it, but
labeled as such, since the difference is worth up to $27.41$ pp of accuracy and $48.66$ pp
of F1 \cite{r33}. Treat benchmark-to-benchmark transfer as insufficient evidence of
generalization: it cost $4.00$ pp in this corpus while benchmark-to-field transfer cost
$14.56$ pp \cite{r49,r55,r61}. Justify each additional sensor against the specific
unimodal baseline it displaces, reporting the ablation under a matched protocol, because the
fusion gain is a function of that baseline and ranged from $+1.00$ to $+14.30$ pp here.
For closed-loop and therapeutic systems \cite{r38,r108}, where a false positive triggers an
intervention rather than a notification, these requirements are not methodological
preferences but safety prerequisites.

\subsection{Limitations}
Four limitations bound these conclusions and should be weighed against them.

\textbf{Small number of within-study contrasts:} The design that gives this review its
internal validity also restricts it: only 4 fusion contrasts and 10 evaluation contrasts met
the matched-protocol requirement. Medians over $k=4$ and $k=10$ are unstable, and the
comparison of the two medians is a comparison of central tendencies in small samples, not a
hypothesis test. We therefore report full ranges throughout, avoid pooled effect-size
estimates and significance testing, and frame the central claim as a difference in order of
magnitude rather than a precise ratio. The personalization estimate in particular rests on
one study reporting two metrics \cite{r33}, and would be materially strengthened or weakened
by two or three further studies reporting both conditions.

\textbf{Corroboration was limited by access, not only by reporting:} A headline outcome was
corroborable for 46.7\% of included studies, and the recovery rates in
Table~\ref{tab:t3} reflect what is obtainable from the openly retrievable record. Studies
behind subscription access may report their protocol and variance fully in the full text.
The rates are therefore a lower bound on the literature's transparency and an accurate
description of what a reader without institutional access can verify. The direction of any
resulting bias is unclear, since there is no reason to expect studies with recoverable
outcomes to differ systematically in performance from those without.

\textbf{Corpus construction:} The corpus was assembled from a pre-specified set of 117 records
rather than from fresh database queries executed for this review, then screened and verified
against strict criteria. This procedure is fully documented and reproducible from
Table~\ref{tab:t1} and Figure~\ref{fig:prisma}, but it is not equivalent to an exhaustive
database search, and studies absent from the source corpus are absent here. The strict
outcome criterion, which excluded 41 records, also narrows coverage relative to reviews that
admit any wearable-sensing application; that narrowing is deliberate but it does reduce
comparability with those reviews.

\textbf{Heterogeneity of metrics and tasks:} Contrasts span accuracy, F1 and AUROC, and
binary through four-class tasks. Differences are reported in percentage points of the
original metric and are never pooled across metrics, but a percentage point of F1 and a
percentage point of AUROC are not interchangeable quantities, and the reader should treat
the aggregate medians as summaries of a heterogeneous set rather than as estimates of a
single underlying parameter.

\section{Conclusion}
Multimodal fusion of wearable and camera-based sensing improves stress and attention
detection, but by less than is generally assumed once cohort, pipeline and validation
protocol are held fixed: a median of $4.67$ percentage points across the protocol-matched
contrasts available in the 2022 to 2026 literature. The same literature shows that evaluation
choices move reported performance roughly twice as far, with personalization alone worth a
median of $38.03$ percentage points and laboratory-to-field transfer costing $13.00$
percentage points of F1. On the field's reference benchmark, classical and transformer models
are separated by under two percentage points, while the validation protocol that produced any
given number is recoverable for only 13.3\% of studies.

The practical consequence is that the next increment of deployable performance is unlikely to
come from another architecture or another sensor. It will come from evaluating honestly:
reporting subject-independent performance as the primary result, validating in the field
rather than only across curated benchmarks, sizing sensor complements against the actual
unimodal baseline, and publishing protocol and variance alongside every point estimate.
Subject-independent generalization on existing benchmarks is unsolved and discriminating, and
is the appropriate target for the next phase of work in this area.
